% =========================================================
% Chapter 1 Locked Refactor: Part A
% =========================================================

\subsection{Algebraic Structures}
\label{sec:algebraic-structures}

\begin{remark*}[Geometric picture]
Before any specific algebra, fix what kind of thing we are studying: a set
of elements together with stipulated ways to combine them, compare them, and
name some of them. The axioms say which stipulations we keep. A group, a
ring, and a field are the same general idea with different stipulations.
\end{remark*}

\begin{definitionbox}{Definition (Algebraic Structure)}
\begin{definition}[Algebraic Structure]
\label{def:algebraic-structure}
An \textbf{algebraic structure} of signature $\mathcal L$ is a pair
\[
  \mathcal A=(A,I)
\]
where $A$ is a nonempty carrier set and $I$ interprets each symbol of
$\mathcal L$ on $A$: each $n$-ary function symbol as an operation
$A^n\to A$, each $n$-ary relation symbol as a subset of $A^n$, and each
constant symbol as an element of $A$. The structure is algebraic when
$\mathcal L$ has no relation symbols beyond equality.
\end{definition}
\end{definitionbox}

\begin{remark*}[Standard quantified statement]
\[
\operatorname{AlgStructure}_{\mathcal L}(\mathcal A)
\Longleftrightarrow
\mathcal A=(A,I),\ A\neq\varnothing,\ \text{and } I
\text{ interprets every symbol of }\mathcal L\text{ on }A.
\]
\end{remark*}

\begin{remark*}[Interpretation]
The signature is the type; the structure is the instance. ``Group'' is a
signature with a binary operation, an identity constant, and an inverse
operation; a particular group interprets these on a carrier.
\end{remark*}

\begin{remark*}[Bourbaki anchor]
Species of structure, \textit{Theory of Sets}, Chapter IV, Section 1.
\end{remark*}

\begin{dependencies}
\begin{itemize}
  \item \hyperref[def:carrier-set]{Carrier Set}
  \item \hyperref[def:signature]{First-Order Signature}
\end{itemize}
\end{dependencies}

\begin{definitionbox}{Definition (First-Order Signature)}
\begin{definition}[First-Order Signature]
\label{def:signature}
A \textbf{first-order signature}
\[
  \mathcal L=(\operatorname{Const}_{\mathcal L},
  \operatorname{Func}_{\mathcal L},
  \operatorname{Rel}_{\mathcal L},\operatorname{arity})
\]
lists constant symbols, function symbols, and relation symbols, with each
function or relation symbol carrying a declared arity.
\end{definition}
\end{definitionbox}

\begin{dependencies}
\begin{itemize}
  \item \hyperref[def:nullary]{Nullary Operation}
  \item \hyperref[def:unary]{Unary Operation}
  \item \hyperref[def:binop]{Binary Operation}
\end{itemize}
\end{dependencies}

\begin{definitionbox}{Definition ($\mathcal L$-Structure)}
\begin{definition}[$\mathcal L$-Structure]
\label{def:l-structure}
Given a signature $\mathcal L$, an \textbf{$\mathcal L$-structure} is a
nonempty carrier set $A$ together with an interpretation of every symbol of
$\mathcal L$ on $A$.
\end{definition}
\end{definitionbox}

\begin{dependencies}
\begin{itemize}
  \item \hyperref[def:signature]{First-Order Signature}
  \item \hyperref[def:carrier-set]{Carrier Set}
\end{itemize}
\end{dependencies}
